\documentclass[11pt]{article}
\usepackage{times}
\usepackage{fullpage}
\usepackage{url}
\usepackage{hyperref}
\usepackage{fancyhdr}
\usepackage{graphicx}
\usepackage{enumitem}
\usepackage{subcaption}
\usepackage{amsmath, amsfonts, amsthm, fouriernc}
\usepackage{IEEEtrantools}
\usepackage{parskip}

\pagestyle{fancy}
\addtolength{\headheight}{2em}
\addtolength{\headsep}{1.5em}
\lhead{\large{ME 2801}}

\usepackage{sectsty}
\sectionfont{\fontsize{12}{8}\selectfont}

\begin{document}

\begin{center}
  \Large{\bf{Assignment: Laplace Transforms}}
\end{center}

% ---------------------------------------------------------------
\section*{Part 1: Forward Laplace Transform}

Using the Laplace transform table, find $F(s) = \mathcal{L}\{f(t)\}$ for each
function below.  Assume $f(t) = 0$ for $t < 0$.

\begin{enumerate}

\item $f(t) = 4e^{-3t}$

\item $f(t) = 2\sin(5t)$

\item $f(t) = 3\cos(2t)$

\item $f(t) = t\,e^{-2t}$

\item $f(t) = 5 - 5e^{-t}$

\end{enumerate}

% ---------------------------------------------------------------
\section*{Part 2: Inverse Laplace Transform by Partial Fractions}

Find $f(t) = \mathcal{L}^{-1}\{F(s)\}$ for each expression using partial fraction
expansion.  Show all steps.

\begin{enumerate}[resume]

\item $F(s) = \dfrac{3}{s(s+3)}$

\item $F(s) = \dfrac{s+1}{(s+2)(s+4)}$

\item $F(s) = \dfrac{2}{s^2 + 4s + 3}$

\item $F(s) = \dfrac{10}{s(s^2 + 4s + 8)}$
  \hfill (Hint: complete the square in the denominator)

\end{enumerate}

% ---------------------------------------------------------------
\section*{Part 3: Solving ODEs via Laplace Transforms}

For each ODE, assume zero initial conditions and a unit step input $r(t) = u(t)$
($R(s) = 1/s$).
(a)~Take the Laplace transform.
(b)~Solve for $C(s)$.
(c)~Find $c(t)$ via partial fractions and inverse transform.

\begin{enumerate}[resume]

\item $\dot{c}(t) + 4c(t) = 8\,r(t)$

\item $\ddot{c}(t) + 5\dot{c}(t) + 6c(t) = 6\,r(t)$

\item $\ddot{c}(t) + 4\dot{c}(t) + 8c(t) = 8\,r(t)$
  \hfill (underdamped: express result in the form
  $c(t) = K[1 - A\,e^{-\sigma t}\cos(\omega_d t - \phi)]$)

\end{enumerate}

% ---------------------------------------------------------------
\section*{Part 4: Pole Locations and Time-Domain Behavior}

\begin{enumerate}[resume]

\item For each transfer function below: identify the poles; state whether the
  system is stable; and describe the qualitative character of the step response
  (e.g., overdamped, underdamped, undamped oscillation, or unstable).

  \begin{enumerate}
    \item $G(s) = \dfrac{1}{(s+1)(s+3)}$
    \item $G(s) = \dfrac{1}{s^2 + 2s + 5}$
    \item $G(s) = \dfrac{1}{s(s+2)}$
    \item $G(s) = \dfrac{1}{(s-1)(s+3)}$
  \end{enumerate}

\item Using MATLAB, create \texttt{tf} objects for each system in the previous
  problem and plot the step responses with \texttt{step()}.
  Confirm that the plots match your qualitative predictions.

\end{enumerate}

\end{document}
